\section{The arithmetic response and its orbit coefficients}
\label{sec:target}

\subsection{The shift equation and the physical requirement}

Define the completed function and its centered logarithmic derivative by
\begin{equation}
 \xi(s)=\tfrac12s(s-1)\pi^{-s/2}\Gamma(s/2)\zeta(s),\qquad
 H(p)=\xi(\tfrac12+p),\qquad m(p)=\frac{H'(p)}{H(p)}.
\end{equation}
The target transfer satisfies
\begin{equation}
 K_\omega(p)=\frac{H(p-\omega)}{H(p+\omega)},\qquad
 \partial_\omega K_\omega=-a_\omega K_\omega,\qquad
 a_\omega(p)=m(p-\omega)+m(p+\omega),\qquad K_0=1.
 \label{eq:shift_equation}
\end{equation}
These identities are initially used on a sufficiently far right,
zero-free half-plane, and then meromorphically. Here $p$ is the Laplace
variable, $\omega$ the arithmetic shift, and $\beta$ the thermal
parameter in \eqref{eq:main_dictionary}; physical time is a separate
variable.

A physical realization must specify its states, preparation, propagation
and readout, and derive $K_\omega$ from them. In particular, on an
ordinary radiation space $L^2(0,L)$ its causal compression
$V_{\omega,L}$ must have the appropriate identity limit and accumulated
norm balance. The desired contraction is
\begin{equation}
 \norm{V_{\omega,L}f}_2^2\le\norm{f}_2^2,
 \qquad V_{0,L}=I.
 \label{eq:radiation_requirement}
\end{equation}
The strong identity limit is the relevant convention; an operator-norm
derivative at zero is not presumed. Cumulative contraction does not
require every shift increment to be contractive, or $a_\omega$ itself
to be a holomorphic positive-real function. Those distinctions and the
unconditional weighted bounds are developed in \cite{SourceNote}.

Writing down the quotient in \eqref{eq:shift_equation} is not a proof
of the full unweighted realization. Its unit boundary modulus, where
defined, follows from the functional equation and conjugation; causal
contractivity for every positive shift is a further issue, related to
the RH-equivalent innerness criterion \cite{Suzuki}. No such conclusion
is inferred here from the thermal coefficient identity.

\subsection{Euler coefficients at unchanged logarithmic labels}

For $\Ree p>1/2+\omega$, factor the target as
\begin{align}
 K_\omega(p)&=A_\omega(p)Z_\omega(p),\qquad
 Z_\omega(p)=\frac{\zeta(p+1/2-\omega)}{\zeta(p+1/2+\omega)},
 \label{eq:factorization}\\
 A_\omega(p)&=\pi^\omega\frac{\Gamma(s_-/2)}{\Gamma(s_+/2)}
             \frac{s_-(s_--1)}{s_+(s_+-1)},\qquad
 s_\pm=p+\tfrac12\pm\omega.
 \label{eq:archimedean}
\end{align}
For one prime, set $a=\ell^{\omega-1/2}$,
$b=\ell^{-\omega-1/2}$ and $z=\ell^{-p}$. Then
\begin{equation}
 \frac{1-bz}{1-az}
 =1+\sum_{r\ge1}(1-\ell^{-2\omega})\ell^{r(\omega-1/2)}z^r.
\end{equation}
Multiplying the absolutely convergent local series gives
\begin{equation}
 Z_\omega(p)=\sum_{n\ge1}c_n(\omega)e^{-p\log n},\qquad
 c_n(\omega)=n^{\omega-1/2}\prod_{\ell\mid n}(1-\ell^{-2\omega}),
 \qquad c_1=1.
 \label{eq:coefficients}
\end{equation}
At $\omega=1/2$, this is $c_n=\varphi_E(n)/n$, with $\varphi_E$
Euler's totient. The orbit terminology refers to the primitive residues
of the modular constant term, not to the closed-geodesic coefficients
of a Selberg trace formula. The labels $\log n$ stay fixed as the
arithmetic weights vary. Whether they are physical delays in a new
model must be derived separately.

The first-prime coefficient and its tangent are
\begin{equation}
 c_2(\omega)=\frac{2^\omega-2^{-\omega}}{\sqrt2},\qquad
 c_2(1/2)=\tfrac12,\qquad c_2'(1/2)=\tfrac32\log2.
 \label{eq:first_prime}
\end{equation}
The full logarithmic derivative also follows in the same convergence
region:
\begin{equation}
 \partial_\omega\log Z_\omega(p)
 =2\sum_{n\ge2}\Lambda(n)n^{-1/2}\cosh(\omega\log n)e^{-p\log n}
 =-a_{\omega,\mathrm{prime}}(p).
 \label{eq:prime_generator}
\end{equation}
The transfer involves all positive integers through multiplicativity;
its logarithmic source involves prime powers through the von Mangoldt
function $\Lambda$. The remaining terms of $-a_\omega$ are
$\partial_\omega\log A_\omega$. A successful physical deformation must
produce those terms with the same parameter.
